\begin{center}
\textsc{Gauss} an \textsc{Olbers}.\quad Göttingen, 14.~Januar 1827.
\end{center}

\par\noindent\hspace*{\parindent}\mbox{.\ .\ .\ .\ .\ .}\quad Ich glaube Ihnen schon früher gemeldet zu haben, dass ich es ganz unthunlich gefunden habe, dasjenige Werk, welches ich über meine Messungen in Zukunft zu geben denke, auch in theoretischer Rücksicht ganz selbstständig zu machen, wenn ich nicht wenigstens einen grossen Theil des Theoretischen vorher anderswo besonders behandle.
Es wird schon voluminös, die Gründe meiner Operationen zu entwickeln, aber, wenn ich mich so ausdrücken darf, die \textls[120]{Gründe} der \textls[120]{Gründe} können nicht in das Werk selbst kommen, ohne es ganz buntscheckig und doch unbefriedigend zu machen.
Ich habe mich daher entschlossen, verschiedene theoretische Materien erst abgesondert in einzelnen Abhandlungen zu entwickeln, wodurch es auch allein möglich wird, diese bedeutenden neuen Capitel der Mathematik mit einem gewissen Grade von Vollständigkeit auszuführen.
Gewissermaassen ist meine Preisschrift über die Transformation der Flächen die erste dieser Abhandlungen; die zweite habe ich vor einigen Monaten der k.\ Societät übergeben als Supplementum theoriae combinationis observationum etc.
Sie enthält die Principien, die als Grundlage der \textls[120]{Ausgleichung} der Beobachtungen angewandt werden müssen, und selbst einige Beispiele von \textsc{Krayenhoffs} und meinen Messungen.\newline
\mbox{.\ .\ .\ .\ .\ .}\par
